\documentclass[11pt,a4paper]{article}

% =============================================================================
% PRACTICAL ASSIGNMENT I - TECHNICAL REPORT TEMPLATE
% Data Stream Mining - NOVA IMS
% Overleaf-ready / pdfLaTeX
%
% STUDENTS: Start by editing ONLY the metadata block below.
% Keep the report structure, margins, font size, and section order unchanged.
% =============================================================================

\usepackage[utf8]{inputenc}
\usepackage[T1]{fontenc}
\usepackage{lmodern}
\usepackage{microtype}
\usepackage[a4paper,top=2.15cm,bottom=2.30cm,left=2.35cm,right=2.35cm,headheight=22pt]{geometry}
\usepackage[table]{xcolor}
\usepackage{graphicx}
\usepackage{booktabs}
\usepackage{tabularx}
\usepackage{array}
\usepackage{ragged2e}
\usepackage{enumitem}
\usepackage{amsmath,amssymb}
\usepackage{tikz}
\usepackage[most]{tcolorbox}
\usepackage{titlesec}
\usepackage{fancyhdr}
\usepackage{lastpage}
\usepackage{hyperref}
\usepackage{bookmark}
\usepackage{pifont}
\usepackage{setspace}
\usepackage{float}
\usepackage{caption}
\usepackage{subcaption}
\usepackage[backend=biber,style=authoryear,maxcitenames=2]{biblatex}
\addbibresource{references.bib}
\usepackage{ifthen}

% -----------------------------------------------------------------------------%
% Packages added by student
% -----------------------------------------------------------------------------%
\usepackage{mathtools}
\usepackage{algorithm}
\usepackage{algpseudocode}

% -----------------------------------------------------------------------------%
% 1. STUDENT METADATA - EDIT THIS BLOCK
% -----------------------------------------------------------------------------%
\newcommand{\AssignmentCode}{PA01}
\newcommand{\StudentA}{Louis Silva}
\newcommand{\StudentAID}{20260613}
\newcommand{\StudentB}{Louis Silva}
\newcommand{\StudentBID}{20260613}
\newcommand{\CityA}{Lisboa}
\newcommand{\CityB}{Porto}
\newcommand{\BusinessScenario}{Urban Walking Tour Operator}
\newcommand{\EventTarget}{T1: Rain Event}
\newcommand{\TargetDefinition}{$y_t = 1$ if $P_t \ge 0.1$ mm; otherwise, $y_t = 0$}
\newcommand{\HistoricalPeriod}{2020-01-01 to 2025-12-31}
\newcommand{\WarmupRule}{Earliest $25\%$ of chronological observations}
\newcommand{\RandomSeed}{42}
\newcommand{\SubmissionDate}{25 October 2026}

% TEMPLATE MODE:
% Keep \draftinstructionstrue while writing to display guidance boxes.
% BEFORE FINAL SUBMISSION change it to: \draftinstructionsfalse
\newif\ifdraftinstructions
% \draftinstructionstrue
\draftinstructionsfalse

% --- Results Table Stuff --- %

\definecolor{LisboaGreen}{HTML}{21A58B}
\definecolor{PortoRed}{HTML}{A5213B}

\newcommand{\splitcell}[2]{%
	\begin{tikzpicture}[baseline=(current bounding box.center)]
		\def\w{1.5}\def\h{0.9}
		% Triangles
		\fill[LisboaGreen!70] (0,\h) -- (\w,\h) -- (0,0) -- cycle;
		\fill[PortoRed!70] (\w,\h) -- (\w,0) -- (0,0) -- cycle;
		% Border and diagonal
		\draw[black] (0,0) rectangle (\w,\h);
		\draw[black] (\w,\h) -- (0,0);
		% Values
		\node[anchor=north west, font=\small, inner sep=2pt] at (0,\h) {#1};
		\node[anchor=south east, font=\small, inner sep=2pt] at (\w,0) {#2};
	\end{tikzpicture}%
}

% -----------------------------------------------------------------------------%
% 2. VISUAL IDENTITY
% -----------------------------------------------------------------------------%
\definecolor{Navy}{HTML}{0B2F4F}
\definecolor{Teal}{HTML}{21A58B}
\definecolor{Mint}{HTML}{9CE3D5}
\definecolor{PaleMint}{HTML}{EEF8F5}
\definecolor{PaleBlue}{HTML}{F2F7FB}
\definecolor{BodyText}{HTML}{16324A}
\definecolor{SecondaryText}{HTML}{526575}
\definecolor{RuleGrey}{HTML}{DFE7EB}
\definecolor{PromptBG}{HTML}{F6F9FB}
\definecolor{WarningBG}{HTML}{FFF8EA}
\definecolor{WarningLine}{HTML}{D28A12}
\definecolor{SoftGrey}{HTML}{F1F3F5}
\definecolor{White}{HTML}{FFFFFF}

\color{BodyText}
\setlength{\parindent}{0pt}
\setlength{\parskip}{0.45em}
\setlist[itemize]{leftmargin=1.6em,itemsep=0.20em,topsep=0.2em}
\setlist[enumerate]{leftmargin=1.8em,itemsep=0.20em,topsep=0.2em}
\renewcommand{\arraystretch}{1.18}
\setstretch{1.03}

\hypersetup{
	colorlinks=true,
	linkcolor=Navy,
	urlcolor=Teal,
	citecolor=Navy,
	pdftitle={Practical Assignment I - Technical Report - \AssignmentCode},
	pdfauthor={\StudentA\ and \StudentB}
}

\pagestyle{fancy}
\fancyhf{}
\fancyhead[L]{\small\color{SecondaryText}\textbf{DATA STREAM MINING} \;|\; Practical Assignment I}
\fancyhead[R]{\small\color{SecondaryText}\AssignmentCode\; | \;\CityA\ vs. \CityB}
\fancyfoot[L]{\small\color{SecondaryText}NOVA IMS}
\fancyfoot[R]{\small\color{SecondaryText}Page \thepage\ of \pageref*{LastPage}}
\renewcommand{\headrulewidth}{0.4pt}
\renewcommand{\headrule}{\hbox to\headwidth{\color{RuleGrey}\leaders\hrule height \headrulewidth\hfill}}
\renewcommand{\footrulewidth}{0pt}

\titleformat{\section}
{\Large\bfseries\color{Navy}}
{\thesection}{0.65em}{}
[\vspace{0.15em}\color{Mint}\titlerule]
\titleformat{\subsection}
{\large\bfseries\color{Teal}}
{\thesubsection}{0.55em}{}
\titlespacing*{\section}{0pt}{1.15em}{0.55em}
\titlespacing*{\subsection}{0pt}{0.95em}{0.35em}

\newtcolorbox{keybox}[1][]{
	enhanced,breakable,colback=PaleMint,colframe=Teal,
	boxrule=0.8pt,leftrule=4pt,arc=2mm,
	left=4mm,right=4mm,top=3mm,bottom=3mm,
	before skip=0.7em,after skip=0.7em,
	fonttitle=\bfseries\color{Navy},#1
}
\newtcolorbox{businessbox}[1][]{
	enhanced,breakable,colback=PaleBlue,colframe=Navy,
	boxrule=0.8pt,leftrule=4pt,arc=2mm,
	left=4mm,right=4mm,top=3mm,bottom=3mm,
	before skip=0.7em,after skip=0.7em,
	fonttitle=\bfseries\color{White},#1
}
\newtcolorbox{warningbox}[1][]{
	enhanced,breakable,colback=WarningBG,colframe=WarningLine,
	boxrule=0.8pt,leftrule=4pt,arc=2mm,
	left=4mm,right=4mm,top=3mm,bottom=3mm,
	before skip=0.7em,after skip=0.7em,
	fonttitle=\bfseries\color{Navy},#1
}
\newtcolorbox{instructionbox}[1][]{
	enhanced,breakable,colback=PromptBG,colframe=RuleGrey,
	boxrule=0.6pt,arc=1.5mm,
	left=3.5mm,right=3.5mm,top=2.5mm,bottom=2.5mm,
	before skip=0.55em,after skip=0.65em,
	fonttitle=\bfseries\color{SecondaryText},#1
}

\newcommand{\instruction}[1]{%
	\ifdraftinstructions
	\begin{instructionbox}[title=What to write here]
		\small\color{SecondaryText}#1
	\end{instructionbox}
	\fi
}

\newcommand{\draftbanner}{%
	\ifdraftinstructions
	\begin{warningbox}[title=Template mode is ON]
		\small This PDF still displays writing guidance. Before final submission, change
		\texttt{\string\draftinstructionstrue} to
		\texttt{\string\draftinstructionsfalse} near the top of \texttt{main.tex} and compile again.
	\end{warningbox}
	\fi
}

\newcolumntype{Y}{>{\RaggedRight\arraybackslash}X}
\newcolumntype{C}[1]{>{\centering\arraybackslash}p{#1}}
\newcolumntype{L}[1]{>{\RaggedRight\arraybackslash}p{#1}}

% Figure helper: the template compiles even before students upload the image.
\newcommand{\reportfigure}[4][0.88\textwidth]{%
	\begin{figure}[H]
		\centering
		\IfFileExists{#2}{%
			\includegraphics[width=#1]{#2}%
		}{%
			\fbox{\parbox[c][4.0cm][c]{0.82\textwidth}{\centering
					\color{SecondaryText}\textbf{Figure placeholder}\\[0.4em]
					Upload \texttt{\detokenize{#2}} to the Overleaf project,\\
					or replace this figure call with your own file.}}
		}
		\caption{#3}
		\label{#4}
	\end{figure}
}

% -----------------------------------------------------------------------------%
% DOCUMENT
% -----------------------------------------------------------------------------%
\begin{document}
	\pagenumbering{gobble}
	
	% COVER PAGE
	\begin{titlepage}
		\pagecolor{Navy}
		\color{white}
		\vspace*{1.25cm}
		{\small\bfseries\color{Mint}\MakeUppercase{Data Stream Mining \;\textbullet\; NOVA IMS}}\\[0.85cm]
		{\Huge\bfseries Practical Assignment I}\\[0.30cm]
		{\fontsize{25}{29}\selectfont\bfseries Technical Report}\\[0.55cm]
		{\Large\color{Mint}Business-Oriented Data Stream Mining}\\[1.05cm]
		
		\begin{tcolorbox}[colback=white!7!Navy,colframe=Mint,boxrule=0.8pt,arc=2mm,left=5mm,right=5mm,top=4mm,bottom=4mm]
			\color{white}
			\begin{tabularx}{\textwidth}{@{}L{4.3cm}Y@{}}
				\textbf{Assignment code} & \AssignmentCode \\
				\textbf{Business scenario} & \BusinessScenario \\
				\textbf{Cities} & \CityA\ vs. \CityB \\
				\textbf{Event target} & \EventTarget \\
				\textbf{Historical period} & \HistoricalPeriod \\
			\end{tabularx}
		\end{tcolorbox}
		
		\vspace{0.75cm}
		\begin{tabularx}{\textwidth}{@{}YY@{}}
			\textbf{Student A} & \textbf{Student B} \\
			\StudentA & \StudentB \\
			\StudentAID & \StudentBID \\
			Primary city: \CityA & Primary city: \CityB \\
		\end{tabularx}
		
		\vfill
		{\large\bfseries Submission date: \SubmissionDate}\\[0.15cm]
		{\normalsize\color{white!75}Expected main-body length: 6--8 pages, excluding cover, references, and optional appendix.}
		
		\ifdraftinstructions
		\vspace{0.55cm}
		{\small\bfseries\color{Mint}TEMPLATE MODE - WRITING GUIDANCE IS CURRENTLY VISIBLE}
		\fi
		\vspace*{0.55cm}
	\end{titlepage}
	\nopagecolor
	\color{BodyText}
	\pagenumbering{arabic}
	\setcounter{page}{1}
	
	\draftbanner

	% =============================================================================
	\section{Business Context and Objective}
	% =============================================================================
	\instruction{Briefly describe the assigned business problem and the operational decision to be supported. Identify the two cities and the event target. Explain why this event matters to the business. End the section with one concise objective or guiding question. Do not make the recommendation yet. Suggested length: about half a page.}
	
	The assigned business is an urban walking tour operator that operates in Lisboa and Porto. Weather conditions, such as rain, can cause a tour to be cancelled, postponed, or rescheduled. As a result, being able to predict these rain events is valuable. In this report, we define a ``rain event'' as the precipitation at time $t$ being greater than or equal to $0.1\mathrm{mm}$.
	
	An important point is that \CityA, for example, could have greater total precipitation than \CityB, or more rain events, but the occurrence of the events in \CityB\ may be more irregular and harder to predict. Hence, the question is not simply ``which city rains the most''. The objective of this study is therefore to determine which of the two cities, \CityA\ and \CityB, offers more predictable conditions for scheduling walking tours with fewer rain interruptions.
	
	% =============================================================================
	\section{Data and Stream Construction}
	% =============================================================================
	\instruction{Describe the Open-Meteo data source, historical period, weather variables, sampling frequency, chronological ordering, warm-up, preprocessing, target construction, and any cleaning decisions. Explain how potential outliers were inspected. Remember: valid extreme weather observations are not automatically errors. State explicitly that future information was not used to construct current inputs. Suggested length: half to one page.}
	
	\subsection{Stream configuration}
	The data collected for this assignment was gathered through the Open-Meteo Historical Weather API \parencite{openmeteo}. \autoref{tab:stream-config} shows the API endpoint used, the historical period of the data, sampling frequency, retrieved variables, warm-up rule, event target, threshold estimation method (in this case it's just an absolute rule), and random seed. Potential outliers were inspected but not removed from the dataset just because they were outliers.
	
	\begin{table}[H]
		\centering
		\caption{Core stream and experimental configuration}
		\label{tab:stream-config}
		\begin{tabularx}{0.94\textwidth}{@{}L{5.2cm}Y@{}}
			\toprule
			\textbf{Item} & \textbf{Configuration used} \\
			\midrule
			Assignment code & \AssignmentCode \\
			API Endpoint URL & \texttt{https://archive-api.open-meteo.com/v1/archive} \\
			City A / City B & \CityA\ / \CityB \\
			Historical period & \HistoricalPeriod \\
			Sampling frequency & Hourly (1 observation per hour) \\
			Variables retrieved & \texttt{temperature\_2m}, \texttt{relative\_humidity\_2m}, \texttt{precipitation}, \texttt{wind\_speed\_10m} \\
			Warm-up rule & \WarmupRule \\
			Event target & \EventTarget \\
			Target definition / threshold & \TargetDefinition \\
			Threshold estimation & Absolute rule \\
			Random seed & \RandomSeed \\
			\bottomrule
		\end{tabularx}
	\end{table}
	
	\subsection{Target and temporal alignment}
	\instruction{Define the binary target mathematically or in precise pseudocode. Explain which observations are used as predictors at time $t$ and why the target-defining value at time $t$ is excluded from the inputs. If a threshold is data-driven, explain that it was estimated from the pooled warm-up only and then frozen.}
	
	As stated in \autoref{tab:stream-config}, the observations are hourly. Hence, observations at $t$ and $t+1$ are spaced one hour apart. For each time step $t$, we define the following observations:
	
	\begin{align}
		T_t &\in \mathbb{R} && \text{as the air temperature at 2m above ground level, measured in } ^\circ\mathrm{C}, \\
		H_t &\in [0,100] && \text{as the relative humidity at 2m above ground level, measured in } \%, \\
		P_t &\in \mathbb{R}_{\geq 0} && \text{as the precipitation accumulated over the preceding hour, measured in mm}, \\
		W_t &\in \mathbb{R}_{\geq 0} && \text{as the wind speed at 10m above ground level, measured in km/h}.
	\end{align}
	
	The binary rain-event target, $y_t \in \{0, 1\}$, is defined as:
	
	\begin{equation}
		\label{eq:rain-event}
		y_t \;\coloneqq\; \begin{cases}
			1 & \text{if} \;\; P_t \:\ge\: 0.1\mathrm{mm}, \\
			0 & \text{if} \;\; P_t \:<\: 0.1\mathrm{mm}.
		\end{cases}
	\end{equation}
	
	To construct the feature vector for predicting $y_t$, we use observations from time $t-1$. This is because the prediction at time $t$ must only use information available before time $t$. Using observations from time $t$ to predict the target at time $t$ would cause temporal leakage. This can result in a deceptively high model performance because the model is using information that would not be available at prediction time.
	
	% =============================================================================
	\section{City Analysis}
	% =============================================================================
	Before we analyse \CityA\ and \CityB, we will define some statistics. Let
	
	\begin{equation}
		W^T_m \;\coloneqq\; \{ y_i : \tau_m \le \tau_i < \tau_{m+1} \},
	\end{equation}
	
	be a tumbling window where $m$ represents a given month, $\tau_m$ is the timestamp at the start of month $m$, $\tau_{m+1}$ is the timestamp at the start of month $m+1$, and $\tau_i$ is the timestamp associated with observation $i$. Then, the percentage of observations where there is a rain event in month $m$ is given by:
	
	\begin{equation}
		\label{eq:R_m}
		R_m \;\coloneqq\; \frac{100}{\left\lvert W^T_m \right\rvert} \sum_{y_i \in W^T_m} y_i, \qquad R_m \in [0, 100].
	\end{equation}
	
	Specifically, $R_m$ denotes the percentage of observations in a given month $m$ that were rain events. In addition, let
	
	\begin{equation}
		W^{(24)}_{t-1} \;\coloneqq\; \{y_{t-24}, y_{t-23}, \dotsc, y_{t-1}\},
	\end{equation}
	
	be a sliding window containing the past $24$ observations. Then, the number of rain-event hours in the past $24$ hours is given by:
	
	\begin{equation}
		\label{eq:C_t^(24)}
		C^{(24)}_t \;\coloneqq \sum_{y_i \in W^{(24)}_{t-1}} y_i, \qquad C^{(24)}_t \in \{0, 1, \dotsc, 24\}.
	\end{equation}
	
	% =============================================================================
	\subsection{City A Analysis: \CityA}
	% =============================================================================
	\instruction{This section is primarily the responsibility of Student A. Characterise the stream for City A before interpreting the models. Include data quality observations, descriptive statistics, target class distribution, temporal behaviour, rolling/window evidence, and at least one analysis that is relevant to the assigned business context. Explain what the evidence means; do not simply describe the shape of a chart. Suggested length: about one page.}

	The \CityA\ stream contains $52{,}608$ hourly observations from 1 January 2020 to 31 December 2025. The notebook checks found no missing values or duplicated rows, no duplicated timestamps, and no non-hourly transitions. With the target defined in \autoref{eq:rain-event}, $6{,}348$ observations ($12.1\%$) are rain events and $46{,}260$ ($87.9\%$) are non-events.
	
	\begin{table}[H]
		\centering
		\caption{Selected descriptive, event, and temporal statistics for \CityA.}
		\label{tab:city-a-profile}
		\begin{tabularx}{0.88\textwidth}{@{}L{6.0cm}C{3.2cm}Y@{}}
			\toprule
			\textbf{Indicator} & \textbf{Value} & \textbf{Why it matters} \\
			\midrule
			Number of observations & $52{,}608$ & Size of the hourly stream \\
			Positive-event rate (\%) & $12.1\%$ & Frequency of observations meeting the rain-event definition \\
			$Q_{0.95}(P_t)$ & $0.3\mathrm{mm}$ & Upper-tail hourly precipitation statistic \\
			$\Pr(H_t \geq 50\% \mid y_t=1)$ & $99.8\%$ & Humidity level among rain-event observations \\
			Median $C^{(24)}_t$ & $0$ hours & Median number of rain-event hours in the preceding 24 observations \\
			$Q_{0.90}(C^{(24)}_t)$ & $10$ hours & Upper-tail count of rain-event hours in the preceding 24 observations \\
			$Q_{0.95}(C^{(24)}_t)$ & $15$ hours & More extreme concentration of rain-event hours in the preceding 24 observations \\
			$Q_{0.90}$ rain-episode duration & $8$ hours & Length of a long but relatively common continuous rain episode \\
			Maximum rain-episode duration & $36$ hours & Longest observed uninterrupted run of rain-event hours \\
			Annual event-rate range & $9.5\%$--$15.6\%$ & Observed variation across calendar years (2020--2025) \\
			\bottomrule
		\end{tabularx}
	\end{table}

	The observed annual rain-event rates ranged from $9.53\%$ in 2023 to $15.57\%$ in 2025. By season, the event rate was lowest in summer in each year, ranging from $1.90\%$ to $4.89\%$. The peak hourly event rate occurred at 14:00 in five of the six years, and at 13:00 in 2022; aggregated across all years, 14:00 had a $20.30\%$ event rate compared with the overall hourly rate of $12.07\%$. Thursday was the highest-risk weekday at $12.98\%$, only modestly above the overall rate, so hour-of-day variation is the clearer observed scheduling signal. These are historical associations, not guarantees about future weather.

	Among rain-event observations, $99.8\%$ had relative humidity of at least $50\%$. For the preceding-24-hour count, $C^{(24)}_t$, the median was $0$, the 90th percentile was $10$, and the 95th percentile was $15$ (out of a possible maximum of 24). This indicates that the typical rolling window is event-free, while the upper tail includes substantial clustering of rain-event hours. Rain-event episodes had a median duration of 2 hours, a 90th-percentile duration of 8 hours, and a maximum duration of 36 hours, evidence that longer wet spells can pose sustained interruption risk even though they are not typical.
	
	\reportfigure{figures/Lisboa_monthly_rain_event_rate.pdf}{Monthly percentage of observations classified as rain events in \CityA\ from \HistoricalPeriod.}{fig:lisboa-monthly-rain-event-rate}
	
	\reportfigure{figures/Lisboa_distribution_of_rainy_hours.pdf}{Distribution of rain-event hours in the preceding 24 observations, $C^{(24)}_t$, for \CityA. The vertical axis uses a logarithmic scale.}{fig:lisboa-distribution-of-rainy-hours}
	
	
	% =============================================================================
	\subsection{City B Analysis: \CityB}
	% =============================================================================
	\instruction{This section is primarily the responsibility of Student B. Use the same analytical logic applied to City A so that the subsequent comparison is fair. Characterise data quality, descriptive behaviour, target distribution, temporal behaviour, and rolling/window evidence. Explain the operational meaning of the evidence. Suggested length: about one page.}
	
	In the \CityB\ stream, there were no missing values, no duplicates, and each observation was exactly one hour after its preceding observation. Furthermore, the gathered weather variables $T_t,\,H_t,\,P_t$ and $W_t$ were all within their correct domains. As a result, no cleaning was necessary. The binary target distribution for $y_t \in \{0,1\}$, defined in \autoref{eq:rain-event}, is:
	
	\begin{align}
		\Pr(y_t = 1) \;&\approx\; 0.204, \\
		\Pr(y_t = 0) \;&\approx\; 0.796.
	\end{align}
	
	This gives a positive-event rate of $20.4\%$, as shown in \autoref{tab:city-b-profile}. Equivalently, this means $20.4\%$ of the observations had at least $0.1\mathrm{mm}$ of precipitation. Interestingly, we have that $Q_{0.95} (P_t) = 0.9\mathrm{mm}$ and $\mathrm{std}(P_t) = 0.67\mathrm{mm}$, but $\max(P_t) = 21.5\mathrm{mm}$. This means that $95\%$ of observations had a precipitation of $0.9\mathrm{mm}$ or below, whilst the maximum indicates the presence of rare extreme precipitation values.
	
	The distribution of $H_t$ is negatively skewed, with a skewness of $-0.78$ and median of $81\%$. An interesting statistic is that among the observations classified as rain events,
	
	\begin{equation}
		\Pr(H_t \geq 50\% \mid y_t = 1) \;\approx\; 0.999.
	\end{equation}
	
	Thus, almost all rain events occurred when the relative humidity was at least $50\%$. The distribution of $T_t$ is relatively symmetric, with the mean and median being approximately equal to $15.5\,^\circ\mathrm{C}$, and the standard deviation being $4.8\,^\circ\mathrm{C}$. The distribution of $W_t$ is positively skewed, with a skewness of $0.86$ and median of $13\,\mathrm{km/h}$.
	
	\autoref{fig:porto-monthly-rain-event-rate} shows the monthly rain-event evolution by plotting $R_m$ (defined in \autoref{eq:R_m}) for all $m$. Intuitively, the rain-event rate tends to peak in winter, and dip in summer. \autoref{fig:porto-distribution-of-rainy-hours} shows the distribution of rainy hours within a 24-hour sliding window by plotting $C^{(24)}_t$. As expected, most of the observation mass is on zero rain-event hours in the preceding 24 hours. However, after a steep decline, the distribution slowly tapers off. This suggests that most of the time it doesn't rain, but when it does, the length of the rainy spell is uncertain.
	
	\begin{table}[H]
		\centering
		\caption{Selected descriptive and event statistics for \CityB.}
		\label{tab:city-b-profile}
		\begin{tabularx}{0.88\textwidth}{@{}L{6.0cm}C{3.2cm}Y@{}}
			\toprule
			\textbf{Indicator} & \textbf{Value} & \textbf{Why it matters} \\
			\midrule
			Number of observations & $52608$ & Size of analysed stream \\
			Positive-event rate (\%) & $20.4\%$ & Frequency of rain-related interruption \\
			$Q_{0.95}(P_t)$ & $0.9\mathrm{mm}$ & Characterises upper-tail hourly precipitation \\
			$Q_{0.90} \left(C^{(24)}_t \right)$ & 17 & Characterises upper-tail concentration of rain events within a 24-hour sliding window \\
			$\mathrm{median} \left(C^{(24)}_t \right)$ & $0$ & Characterises the typical amount of rain events within a 24-hour sliding window \\
			\bottomrule
		\end{tabularx}
	\end{table}
	
	\reportfigure{figures/Porto_monthly_rain_event_rate.pdf}{Plotting the monthly rain-event rate $R_m$ for \CityB\ from \HistoricalPeriod.}{fig:porto-monthly-rain-event-rate}
	
	\reportfigure{figures/Porto_distribution_of_rainy_hours.pdf}{Plotting the distribution of rainy hours in the past $24$ hours, $C^{(24)}_t$, for \CityB\ from \HistoricalPeriod. The $y$-axis is plotted on a log scale.}{fig:porto-distribution-of-rainy-hours}
	
	
	% =============================================================================
	\section{Comparative Data Analysis}
	\label{sec:comparative-data-analysis}
	% =============================================================================
	\instruction{Compare the two cities before discussing classifier scores. Focus on differences that matter to the assigned business question: event frequency, variability, persistence, temporal patterns, operational thresholds, or other justified indicators. The comparison must use compatible definitions for both cities. This section is joint work. Suggested length: about one page.}
	
	To compare the business-relevant differences in the two cities, we will define two additional statistics. Let $B_i$ denote an event where for some observation indices $s_i$ and $e_i$, where $s_i \le e_i$, all observations $y_{s_i} = y_{s_i + 1} = \dotsc = y_{e_i} = 1$, and $y_{s_i - 1} = 0$ and $y_{e_i + 1} = 0$. Therefore, $B_i$ describes the $i$-th maximal sequence of observations where there was a rain event, informally known as a ``rainy period''. Formally,
	
	\begin{equation}
		B_i \;\coloneqq\; \{s_i,\,s_i+1,\,\dotsc,\,e_i\}
		\quad \text{such that} \quad
		\begin{cases}
			y_\tau = 1, \; \forall \tau \in B_i, \text{ and} \\
			y_{s_i-1} = y_{e_i+1} = 0.
		\end{cases}
	\end{equation}
	
	Therefore, if we observed $b \in \mathbb{N}$ total rainy periods, then
	
	\begin{equation}
		\mathrm{median} \left(\left(\lvert B_i \rvert\right)_{i=1}^b \right)
	\end{equation}
	
	denotes the median length of these rainy periods. Similarly, we can define the $i$-th ``dry period'' as
	
	\begin{equation}
		D_i \;\coloneqq\; \{s_i,\,s_i+1,\,\dotsc,\,e_i\}
		\quad \text{such that} \quad
		\begin{cases}
			y_\tau = 0, \; \forall \tau \in D_i, \text{ and} \\
			y_{s_i-1} = y_{e_i+1} = 1.
		\end{cases}
	\end{equation}
	
	Consequently, if we observed $d \in \mathbb{N}$ total dry periods, then
	
	\begin{equation}
		\mathrm{median} \left(\left(\lvert D_i \rvert\right)_{i=1}^d \right)
	\end{equation}
	
	denotes the median length of these periods. \autoref{tab:business-comparison} shows these two statistics for both cities. As shown, \CityA\ has a smaller median length of rainy periods, and a larger median length of dry periods than \CityB. This suggests that \CityB\ simultaneously has shorter dry periods and longer rainy periods. Furthermore, \CityB\ also has a higher variability in its monthly rain-event rate, denoted by $\mathrm{std}(R_m)$. This can also be seen visually in panel (a) of \autoref{fig:city-comparison}: in the winter months, the rain-event rate for \CityB\ spikes much higher than \CityA.
	
	In addition, the mean number of rain-event hours within a 24-hour sliding window, given that there is at least one rain-event hour, is higher for \CityB. This suggests that, conditional on a 24-hour window containing at least one rain-event hour, rain-event hours are more concentrated within such windows in \CityB. This can be seen visually in panel (b) of \autoref{fig:city-comparison}: beyond approximately eight rain-event hours per 24-hour window, the proportions for \CityA\ declines much more sharply than those for \CityB. Furthermore, \CityB\ retains relatively larger proportions at higher values of $C^{(24)}_t$.
	
	For walking tours, shorter dry periods and longer rainy periods imply fewer uninterrupted windows in which the tours can be scheduled without rain-related disruption.
	
	\begin{table}[H]
		\centering
		\caption{Business-relevant comparison before modelling.}
		\label{tab:business-comparison}
		\renewcommand{\arraystretch}{1.3}
		\begin{tabularx}{0.92\textwidth}{@{}Y C{3.0cm} C{3.0cm}@{}}
			\toprule
			\textbf{Indicator} & \textbf{\CityA} & \textbf{\CityB} \\
			\midrule
			Positive-event rate (\%) & $12.1$ & $20.4$ \\
			$\mathrm{std}(R_m)$ (\%) & $9.3$ & $12.7$ \\
			$\mathbb{E}\left[C^{(24)}_t \mid C^{(24)}_t > 0\right]$ (hours) & $7.1$ & $10.3$ \\
			$\mathrm{median} \left(\left(\lvert B_i \rvert\right)_{i=1}^b \right)$ & $2$ & $3$ \\
			$\mathrm{median} \left(\left(\lvert D_i \rvert\right)_{i=1}^d \right)$ & $4$ & $3$ \\
			\bottomrule
		\end{tabularx}
	\end{table}
	
	\reportfigure{figures/city\_comparison.pdf}{Comparison of monthly rain-event rates and the distributions of rain-event hours in preceding 24-hour sliding windows for \CityA\ and \CityB.}{fig:city-comparison}
	
	
	% =============================================================================
	\section{Data Stream Mining Experiment}
	% =============================================================================
	\instruction{Describe the two required feature representations, the five mandatory learners, any optional model used, model parameters, the prequential protocol, and the evaluation metrics. Explain how the same chronological history was preserved within each city. If you investigate one model parameter, state the values compared and the criterion used to interpret the trade-off. Suggested length: one to one-and-a-half pages.}
	
	\subsection{Feature representations}
	For Representation A, we define the feature vector $X^A_t$ using the lagged weather variables for temperature, humidity, precipitation, and wind:
	
	\begin{equation}
		X_t^A \;\coloneqq\; \left(T_{t-1},\,H_{t-1},\,P_{t-1},\,W_{t-1}\right).
	\end{equation}
	
	For Representation B, we first need to define some features. Let
	
	\begin{align}
		\Delta P_t &\;\coloneqq\; P_{t-1} \:-\: P_{t-2}, \\
		\Delta H_t &\;\coloneqq\; H_{t-1} \:-\: H_{t-2},
	\end{align}
	
	be the delta features for $P_t$ and $H_t$ respectively. Next, let
	
	\begin{align}
		\bar{P}_t^{(6)} &\;\coloneqq\; \frac{1}{6} \sum_{i=1}^{6} P_{t-i}, \\
		\bar{H}_t^{(6)} &\;\coloneqq\; \frac{1}{6} \sum_{i=1}^{6} H_{t-i},
	\end{align}
	
	be the rolling means for $P_t$ and $H_t$ in a 6-hour sliding window. Then, we define the rolling 6-hour standard deviation for $P_t$ and $H_t$ as:
	
	\begin{align}
		\tilde{P}_t^{(6)} &\;\coloneqq\; \sqrt{\frac{1}{5} \sum_{i=1}^{6} \left(P_{t-i} \:-\: \bar{P}_t^{(6)}\right)^2}, \\
		\tilde{H}_t^{(6)} &\;\coloneqq\; \sqrt{\frac{1}{5} \sum_{i=1}^{6} \left(H_{t-i} \:-\: \bar{H}_t^{(6)}\right)^2}.
	\end{align}
	
	In addition to these new features, we also use the amount of rain events within a 24-hour sliding window, $C_t^{(24)}$ (defined in \autoref{eq:C_t^(24)}), and $X_t^A$, to construct the feature vector for Representation B:
	
	\begin{equation}
		X_t^B \;\coloneqq\; \left(X_{t}^A,\, \Delta P_t,\, \Delta H_t,\, \bar{P}_t^{(6)},\, \bar{H}_t^{(6)},\, \tilde{P}_t^{(6)},\, \tilde{H}_t^{(6)},\, C^{(24)}_{t}\right).
	\end{equation}
	
	\subsection{Model configuration}
	
	\autoref{tab:model-config} shows the five models used and their selected parameters.
	
	\begin{table}[H]
		\centering
		\frenchspacing
		\caption{Models and main configurations used in the experiment.}
		\label{tab:model-config}
		\begin{tabularx}{0.96\textwidth}{@{}L{4.7cm}Y@{}}
			\toprule
			\textbf{Model} & \textbf{Main configuration / parameters} \\
			\midrule
			Gaussian Naive Bayes & Default \\
			Online Logistic Regression & \texttt{l2: 0}, \texttt{standardize: true}, \texttt{optimizer: Stochastic Gradient Descent} with \texttt{learning-rate: 0.05},  \\
			Hoeffding Tree & \texttt{grace\_period: 50}, \texttt{delta: 0.001}, \texttt{tau: 0.10}, \texttt{leaf\_prediction: "mc"} \\
			Online Bagging & \texttt{base-learner: Hoeffding Tree} with the same parameters as defined above, \texttt{n\_models: 5}, \texttt{seed: 42} \\
			Adaptive Random Forest & \texttt{n\_models: 5}, \texttt{max\_features: "sqrt"}, \texttt{lambda\_value: 6}, \texttt{grace\_period: 30}, \texttt{delta: 0.001}, \texttt{tau: 0.01}, \texttt{leaf\_prediction: "mc"}, \texttt{seed: 42} \\
			\bottomrule
		\end{tabularx}
	\end{table}
	
	\subsection{Evaluation protocol}	
	\instruction{State explicitly that each prediction was evaluated before the current label was learned. Report Accuracy, Macro-F1, Precision, Recall and the confusion matrix, together with prediction and update cost per event. Explain which class is treated as the positive event class.}
	
	Algorithm \autoref{alg:prequential-evaluation-protocol} shows the evaluation protocol defined for this experiment. As shown, each prediction is evaluated before the current label is learnt. As defined in \autoref{eq:rain-event}, the positive event class, $y_t = 1$, denotes that at time $t$, $P_t \ge 0.1\mathrm{mm}$.
	
	\begin{algorithm}[H]
		\caption{Prequential evaluation protocol used in the \texttt{ExperimentResults.run()} function in \texttt{src/results.py}}
		\label{alg:prequential-evaluation-protocol}
		\begin{algorithmic}[1]
			\ForAll{cities $c \in \{\text{\CityA}, \text{\CityB}\}$}
				\ForAll{representations $r \in \{A,B\}$}
					\ForAll{models $m$}
						\State Initialise a fresh instance of model $m$
						\State Warm-up $m$ using the earliest $25\%$ of chronological observations
						\ForAll{observations $t$ in chronological order}
							\State \textbf{Predict:} Obtain $\hat{y}_t$ using $X_t^r$ and record prediction time
							\State \textbf{Evaluate:} Update metrics and confusion matrix using $(y_t, \hat{y}_t)$
							\State \textbf{Learn:} Update $m$ using $(X^r_t, y_t)$ and record update time
						\EndFor
						\State Save final metrics, confusion matrix, and mean processing times
					\EndFor
				\EndFor
			\EndFor
		\end{algorithmic}
	\end{algorithm}
	
	
	% =============================================================================
	\section{Results and Discussion}
	% =============================================================================
	\instruction{Present only the evidence needed to answer the assignment. Interpret results rather than ranking models. Discuss class imbalance, event detection, differences between Representation A and Representation B, model behaviour over time, and computational cost. Connect model evidence back to the stream characteristics identified earlier. At least two analyses or visualisations in the complete report must be original group-designed analyses rather than direct reproductions of the starter-code figures. Suggested length: one to one-and-a-half pages.}
	
	\begin{table}[htbp]
		\centering
		\caption
		{
			Predictive and computational performance for Representation A.
			Predictive metrics are reported as percentages.
			Each cell reports
			\colorbox{LisboaGreen!70}{\scriptsize Lisboa}
			in the top-left triangle and
			\colorbox{PortoRed!70}{\scriptsize Porto}
			in the bottom-right triangle.
		}
		\label{tab:representation-a-results}
		
		\resizebox{\textwidth}{!}{%
			\begin{tabularx}{0.98\textwidth}{@{}Y C{1.55cm} C{1.55cm} C{1.55cm} C{1.55cm} C{1.55cm} C{1.55cm}@{}}
				\toprule
				\textbf{Model}
				& \textbf{Acc.}
				& \textbf{Macro-F1}
				& \textbf{Prec.}
				& \textbf{Recall}
				& \textbf{Pred. ms/event}
				& \textbf{Update ms/event} \\
				\midrule
				GaussianNB & \splitcell{92.1}{89.2} & \splitcell{75.1}{79.4} & \splitcell{84.8}{90.8} & \splitcell{40.2}{50.8} & \splitcell{0.052}{0.052} & \splitcell{0.002}{0.002} \\[4pt]
				Online Logistic & \splitcell{92.5}{90.4} & \splitcell{77.4}{82.7} & \splitcell{83.1}{88.2} & \splitcell{45.7}{59.7} & \splitcell{0.002}{0.002} & \splitcell{0.005}{0.005} \\[4pt]
				Hoeffding Tree & \splitcell{93.4}{92.2} & \splitcell{83.6}{87.8} & \splitcell{74.4}{80.4} & \splitcell{67.8}{80.6} & \splitcell{0.005}{0.006} & \splitcell{0.009}{0.009} \\[4pt]
				Online Bagging & \splitcell{93.5}{92.2} & \splitcell{84.2}{87.8} & \splitcell{73.1}{80.5} & \splitcell{71.0}{80.5} & \splitcell{0.031}{0.031} & \splitcell{0.049}{0.050} \\[4pt]
				Adaptive Random \allowbreak Forest & \splitcell{92.6}{90.4} & \splitcell{79.9}{84.4} & \splitcell{75.1}{78.8} & \splitcell{55.7}{71.0} & \splitcell{0.027}{0.028} & \splitcell{0.120}{0.140} \\
				\bottomrule
			\end{tabularx}
		}
	\end{table}
	
	\begin{table}[htbp]
		\centering
		\caption
		{
			Predictive and computational performance for Representation B.
			Predictive metrics are reported as percentages.
			Each cell reports
			\colorbox{LisboaGreen!70}{\scriptsize Lisboa}
			in the top-left triangle and
			\colorbox{PortoRed!70}{\scriptsize Porto}
			in the bottom-right triangle.
		}
		\label{tab:representation-b-results}
		
		\resizebox{\textwidth}{!}{%
			\begin{tabularx}{0.98\textwidth}{@{}Y C{1.55cm} C{1.55cm} C{1.55cm} C{1.55cm} C{1.55cm} C{1.55cm}@{}}
				\toprule
				\textbf{Model}
				& \textbf{Acc.}
				& \textbf{Macro-F1}
				& \textbf{Prec.}
				& \textbf{Recall}
				& \textbf{Pred. ms/event}
				& \textbf{Update ms/event} \\
				\midrule
				GaussianNB & \splitcell{91.5}{89.7} & \splitcell{77.1}{82.4} & \splitcell{69.4}{81.2} & \splitcell{51.2}{63.1} & \splitcell{0.062}{0.062} & \splitcell{0.004}{0.004} \\[4pt]
				Online Logistic & \splitcell{92.7}{90.6} & \splitcell{79.3}{83.9} & \splitcell{79.3}{83.2} & \splitcell{51.7}{65.9} & \splitcell{0.003}{0.003} & \splitcell{0.009}{0.009} \\[4pt]
				Hoeffding Tree & \splitcell{93.2}{92.0} & \splitcell{83.8}{87.5} & \splitcell{71.7}{79.6} & \splitcell{71.0}{80.4} & \splitcell{0.006}{0.006} & \splitcell{0.016}{0.017} \\[4pt]
				Online Bagging & \splitcell{93.5}{92.2} & \splitcell{84.1}{87.8} & \splitcell{73.1}{80.1} & \splitcell{70.9}{80.8} & \splitcell{0.033}{0.033} & \splitcell{0.082}{0.086} \\[4pt]
				Adaptive Random \allowbreak Forest & \splitcell{92.2}{90.3} & \splitcell{79.0}{84.3} & \splitcell{73.1}{77.8} & \splitcell{54.3}{71.5} & \splitcell{0.027}{0.027} & \splitcell{0.144}{0.165} \\
				\bottomrule
			\end{tabularx}
		}
	\end{table}
	
	\autoref{tab:representation-a-results} and \autoref{tab:representation-b-results} show the accuracy, macro-F1, precision, and recall for each model, for Representation A and B respectively, followed by the prediction and update time in milliseconds for each event.
	
	For Gaussian naive Bayes and online logistic regression, Representation B provided an improvement over Representation A in macro-F1, and a large improvement in recall. The improvement in recall was accompanied by a decline in precision for both models in both cities. Consequently, these two models clearly benefited from the temporal features added in Representation B in terms of macro-F1 and recall, although precision declined. 
	
	For Hoeffding trees, Representation B provided a $0.2$ percentage point improvement in macro-F1 for \CityA, but a $0.3$ percentage point deterioration for \CityB. Similarly, recall improved under Representation B for \CityA, but declined for \CityB. Furthermore, precision declined for both cities under Representation B. For online bagging, macro-F1 declined slightly for \CityA\ but remained unchanged for \CityB. For adaptive random forest, macro-F1 declined slightly for both cities under Representation B.
	
	Across all five models and both representations, \CityB\ achieves higher macro-F1, precision, and recall than \CityA. Recall is especially important because it measures the proportion of times the model predicted it was going to rain, given that it actually rained. Conversely, across all five models and both representations, \CityA\ achieves higher accuracy than \CityB, which may partly reflect its larger proportion of dry observations. Interestingly, in \autoref{sec:comparative-data-analysis} we discussed how \CityB\ had a higher monthly rain-event rate variability, and a larger upper tail in its distribution of rain-event hours in the preceding 24 hours. One possible contributing factor is that \CityB\ has a higher positive-event rate and more balanced target distribution than \CityA.
	
	\autoref{fig:model-performance-over-time} shows two panels, one for \CityA\ and one for \CityB, where the rolling macro-F1 score over a 720-observation (approximately 30-day) window was plotted over time for the online bagging and Hoeffding tree models under Representation A. For \CityB, the Hoeffding tree model almost perfectly follows the online bagging model. For \CityA\, however, their respective performances are more varied. For both models in \CityB, the rolling macro-F1 occasionally declines sharply before recovering. In \CityA, these declines appear more frequent and variable. The sharp declines to approximately $50\%$ occur during windows with very few rain events, where the models failed to detect any of those events despite achieving an accuracy close to $100\%$. This illustrates how high accuracy can conceal missed rain events in imbalanced streams.
	
	For all models, the update time was larger with Representation B compared to A. Similarly, prediction time generally increased with Representation B, although it remained unchanged for Hoeffding trees in \CityB\ and adaptive random forest in \CityA, and declined by $0.001$ milliseconds for adaptive random forest in \CityB. The Hoeffding tree model is considerably faster than online bagging and adaptive random forest, whilst achieving comparable macro-F1 and recall to online bagging, and generally higher macro-F1 and recall than adaptive random forest.
	
	In conclusion, the strongest model in terms of macro-F1 was online bagging. It achieved the highest macro-F1 in both cities under both representations, except that it tied with Hoeffding trees in \CityB\ under Representation A. Although the prediction and update times of online bagging were approximately 5 times those of Hoeffding trees, the observations in the data stream arrive hourly, so the additional processing time is unlikely to have an operational impact. Nevertheless, online bagging's small macro-F1 advantage must be weighed against its higher computational cost.
	
	\reportfigure{figures/model\_performance\_over\_time.pdf}{Rolling macro-F1 over a 720-observation (approximately 30-day) window for the online bagging and Hoeffding tree models under Representation A in (a) \CityA\ and (b) \CityB.}{fig:model-performance-over-time}
	
	
	
	\instruction{After the tables/figures, explain the relationships you observed. Examples of useful questions: Did higher Accuracy coincide with higher Recall? Did temporal features improve event detection? Did an ensemble provide a meaningful gain relative to its computational cost? Did one city's class distribution make the task more difficult?}
	
	
	% =============================================================================
	\section{Business and Technical Recommendations}
	% =============================================================================
	\instruction{Provide two separate conclusions. The business recommendation must be supported by descriptive, temporal and operational evidence, not by whichever classifier achieved the highest score. The technical recommendation must be supported by predictive and computational evidence. State the scope of each recommendation carefully. Suggested length: about half a page.}
	
	\begin{businessbox}[title=Business Recommendation]
		\autoref{tab:business-comparison} and \autoref{fig:city-comparison} shows that \CityA\ has a smaller median length of rainy periods, larger median length of dry periods, a lower mean number of rain-event hours in the preceding 24 hours, conditional on at least one rain-event occurring in that window, a smaller monthly rain-event rate variability, and a smaller positive-event rate. Based on these indicators, our recommendation is that \CityA\ is better for scheduling walking tours due to its fewer rain interruptions compared to \CityB.
	\end{businessbox}
	
	\begin{keybox}[title=Technical Recommendation]
		Our technical recommendation for a streaming modelling approach for the continuous prediction of \EventTarget\ is to use the Hoeffding tree model with Representation A, with the parameters defined in \autoref{tab:model-config}. The Hoeffding tree model provides a favourable trade-off between macro-f1, recall, and computational cost. Furthermore, Representation B nearly doubled the update time of the Hoeffding tree, whilst only improving recall by 3.2 percentage points in \CityA\, but reducing it by 0.2 percentage points in \CityB. Online bagging provides a marginally better macro-F1 and recall, but is approximately 5 times slower.
	\end{keybox}
	
	
	% =============================================================================
	\section{Limitations}
	% =============================================================================
	\instruction{Explain what the study cannot establish. Consider the limited set of weather variables, the selected historical period, the simplified binary event definition, external business factors not included in Open-Meteo, computational measurements being environment-specific, and any modelling assumptions. Do not introduce new results here.}
	
	Firstly, the model hyperparameters were not tuned. Secondly, the study only used a very limited set of weather variables. 
	
	The observations only came hourly, and then we constructed a rain-event if there was 0.1mm or more of precipitation in that hour. That 0.1mm of precipitation couldve been concentrated on a much smaller period of time, e.g., 10 minutes, but then we just say that the whole hour is rainy.
	
	The binary rain-event was not based on any empirical evidence. We did not actually measure that 0.1mm or more of precipitation per hour actually produces enough rain to permit cancelling a tour.
	
	The model doesnt take into account extra-ordinary weather events like prolonged storms
	
	Weather is chaotic?
	
	
	% =============================================================================
	\section{AI Assistance Statement}
	% =============================================================================
	\instruction{This statement is mandatory. Adapt one of the formulations below truthfully. Delete the option that does not apply. Generative AI may assist with programming tasks, but it may not perform the interpretation of experimental results, discussion, business recommendation, technical recommendation, or evidence-based conclusions.}
	
	Generative AI tools were used to assist with Python debugging, plotting syntax, LaTeX syntax, and code organisation. They were not used to interpret experimental results or produce the analysis, discussion, business recommendation, technical recommendation, or evidence-based conclusions.
	
	% =============================================================================
	% REFERENCES
	% =============================================================================
	\section*{References}
	\addcontentsline{toc}{section}{References}
	\instruction{Cite the data source, libraries, algorithms/methods, and any external references used. Add BibTeX entries to \texttt{references.bib} and cite them in the report with commands such as \texttt{\string\parencite\{key\}}. Remove references that are not cited.}
	
	\printbibliography[heading=none]
	
	
	% =============================================================================
	% OPTIONAL APPENDIX
	% =============================================================================
	\clearpage
	\appendix
	\section{Optional Appendix}
	\instruction{Use the appendix only for supporting material that is useful for auditability but not essential to the main argument: additional confusion matrices, parameter grids, supplementary plots, or compact implementation notes. The appendix does not compensate for missing interpretation in the main report. Delete this appendix if it is not needed.}
	
	
	% =============================================================================
	% DRAFT-ONLY SELF-CHECK - DISAPPEARS WHEN draftinstructionsfalse
	% =============================================================================
	\ifdraftinstructions
	\clearpage
	\section*{Template Self-Check Before Final Submission}
	\begin{instructionbox}[title=Turn off template mode only after every item below is satisfied]
		\begin{tabularx}{\textwidth}{@{}C{0.8cm}Y@{}}
			\fbox{\rule{0pt}{1.0ex}\rule{1.0ex}{0pt}} & Metadata on the cover are correct: PA code, students, cities, scenario, target and period. \\
			\fbox{\rule{0pt}{1.0ex}\rule{1.0ex}{0pt}} & The report uses chronological data and documents the warm-up. \\
			\fbox{\rule{0pt}{1.0ex}\rule{1.0ex}{0pt}} & No current target-defining value or future observation is used as an input. \\
			\fbox{\rule{0pt}{1.0ex}\rule{1.0ex}{0pt}} & Thresholds derived from data use pooled warm-up information only and remain frozen during evaluation. \\
			\fbox{\rule{0pt}{1.0ex}\rule{1.0ex}{0pt}} & City A and City B use compatible definitions and evaluation logic. \\
			\fbox{\rule{0pt}{1.0ex}\rule{1.0ex}{0pt}} & The five mandatory models are evaluated with PREDICT $\rightarrow$ EVALUATE $\rightarrow$ LEARN. \\
			\fbox{\rule{0pt}{1.0ex}\rule{1.0ex}{0pt}} & Accuracy is interpreted together with Macro-F1, Precision, Recall and the confusion matrix. \\
			\fbox{\rule{0pt}{1.0ex}\rule{1.0ex}{0pt}} & At least two analyses/visualisations are group-designed rather than direct starter-code reproductions. \\
			\fbox{\rule{0pt}{1.0ex}\rule{1.0ex}{0pt}} & Business and Technical Recommendations are separate and evidence-based. \\
			\fbox{\rule{0pt}{1.0ex}\rule{1.0ex}{0pt}} & Every figure/table included is discussed in the report and has a meaningful caption. \\
			\fbox{\rule{0pt}{1.0ex}\rule{1.0ex}{0pt}} & The AI Assistance Statement is truthful and complete. \\
			\fbox{\rule{0pt}{1.0ex}\rule{1.0ex}{0pt}} & The main body is approximately 6--8 pages after template guidance is hidden. \\
			\fbox{\rule{0pt}{1.0ex}\rule{1.0ex}{0pt}} & \texttt{\string\draftinstructionsfalse} is set before exporting the final PDF. \\
		\end{tabularx}
	\end{instructionbox}
	\fi
	
\end{document}
